\section{Introduction}
\label{sec:intro}

Cover bidding---submitting bids designed to lose---is a prevalent
operational mechanism of bid-rigging cartels
\citep{oecd2012fighting}, yet its internal
structure remains poorly understood. How do cartels calibrate cover
bids? Does coordination tighten or widen the bid distribution? And
what does the answer imply for the screens that competition authorities
use to detect collusion? This paper develops a structural model of cover-bidder deployment
and estimates it using an externally constructed frequent-loser (FL)
partition. The estimation yields what we term a \emph{dispersion paradox}:
conditional on the FL labels, a coordinated cover-bidding regime
fits the bid data much better than the complementary alternative,
and under that regime, FL-labeled bids are \emph{less}, not more,
dispersed than genuine bids---the opposite of what the dominant
detection screens assume.

The model specifies two regimes that differ in coordination capacity.
Under \emph{complementary} cover bidding (Regime~1), cover bidders
receive loose instructions (``bid above the winner'') and spread bids
across a wide range, producing high dispersion. Under
\emph{coordinated} cover bidding (Regime~2), cover bidders calibrate
bids near a focal point just above the winning price, producing low
dispersion. The cartel chooses the number of cover bidders, their bid
distribution, and which markets to target as a function of genuine
competition, detection probability, and the institutional rules of the
procurement system. The model generates five testable predictions that
we take to data from 40~million bids in Brazilian public procurement
(2009--2019).

Maximum likelihood estimation on bid-level data selects Regime~2:
cover-bid dispersion is 28\% below genuine-bid dispersion
($\hat{\sigma}_c / \hat{\sigma}_g = 0.72$;
$\Delta$BIC~$= -91{,}473$). Cover bids cluster 0.83 log points above the winning price
($\hat{\epsilon} = 0.83$, approximately $e^{0.83} - 1 \approx 129\%$
in levels), high enough to guarantee losing yet calibrated to
simulate competitive bidding. The
structural estimates are conditional on a partition of firms into
\emph{frequent losers} (FL)---2,735 firms that never win a single
tender yet participate in abnormally many procurement
processes---and genuine bidders. The FL classification is developed
and validated in a companion paper
\citep{genicolomartins2026screening}, where it achieves
AUC~$= 0.94$ against CADE convictions in a temporal holdout (98
confirmed co-bidders among 16,843 always-losers). Here we take the
FL partition as given and estimate a structural mixture model on
189,381 FL-labeled losing bids benchmarked against 23.2~million
genuine losing bids. The structural results inherit the noise and
scope of the FL classification: the model characterizes the
bid-generating process conditional on this partition, not the
identity of true cartel members.

\subsection{The dispersion paradox and its implications}

The dispersion pattern---that a coordinated regime fits FL-labeled
bids better and implies lower dispersion---has three implications.

\emph{First}, it provides one structural interpretation of a gap in
the detection literature. Variance-based screens
\citep{imhof2019detecting} achieve AUC~$= 0.79$ against CADE
convictions, while the participation-based FL screen achieves
AUC~$= 0.94$. Under the model, variance screens are designed for
the complementary regime but may be misaligned with the coordinated
regime, whose bid patterns can resemble competitive tenders. The
participation footprint---never winning yet bidding repeatedly---is
invariant to the bid distribution, making it robust to both regimes.

\emph{Second}, it provides a structural interpretation for the
\citet{bajari2003deciding} exchangeability test. FL bid residuals
violate exchangeability and exhibit excess pairwise correlation,
but these violations \emph{reverse} under tender fixed effects---the
FL pairwise product drops below the non-FL product. Under the
coordinated regime, this is expected: cover bids share a
focal-point component ($b^* + \epsilon$) that tender fixed effects
absorb, leaving only residual dispersion, which under Regime~2 is
low. The reversal is a directional prediction generated by the
focal-point model (Section~\ref{sec:spec_tests}).

\emph{Third}, it yields testable market-selection predictions.
Because cover bidding is most valuable in competitive markets (low
HHI), the model predicts that FL activity concentrates in
competitive rather than concentrated markets. The data are
consistent with this, and the calibrated complementarity elasticity
$\hat{\gamma} = 0.694$ indicates that, conditional on the FL
partition, cover-bidder deployment scales nearly proportionally with
genuine competition---roughly, one additional genuine bidder is
associated with 0.7 additional cover bidders---a relationship that
holds both across and within PBUs
(Section~\ref{sec:results_calibration}).

\subsection{Contributions}

We make three contributions.

\emph{First}, we develop a structural model of cover-bidder
deployment that distinguishes complementary from coordinated
cover-bidding regimes. The model formalizes the cartel's trade-off
between detection risk, cover-bidding cost, and the institutional
constraint of minimum-bidder rules, and generates predictions for the
number, distribution, and market selection of cover bids.
Conditional on the FL partition, Regime~2 (coordinated) fits the
data much better ($\Delta$BIC~$= -91{,}473$): in this setting,
\emph{low} dispersion, not high, characterizes FL-labeled bids.

\emph{Second}, we examine the model's consistency with five
out-of-sample moments not targeted in estimation:
(i)~the convite--preg\~ao deployment asymmetry;
(ii)~within-tender bid dispersion (FL bids less dispersed in
82.6\% of mixed tenders); (iii)~the $m$--$\theta$ comparative
static across PBU size quartiles; (iv)~the Bajari--Ye
exchangeability and conditional-independence violations; and
(v)~the market-selection gradient across winner HHI. All five are directionally consistent with the model's predictions,
though they differ in how cleanly they discriminate the mechanism
from alternatives (Section~\ref{sec:spec_tests}).

\paragraph{What the paper does and does not claim.}
The structural framework organizes three features that reduced-form
analysis alone does not: (i)~a two-regime taxonomy that gives the
dispersion pattern its theoretical interpretation; (ii)~the
Bajari--Ye tender-FE reversal, a directional prediction generated by
the focal-point model; and (iii)~the calibrated complementarity
elasticity ($\hat{\gamma} = 0.694$) that quantifies how cover-bidder
deployment covaries with genuine competition. The paper does
\emph{not} identify true cartel membership (the FL labels are
statistical, not forensic), does not causally identify the
FL--price association (which reflects both treatment and selection),
and does not establish a universal property of cover bidding across
procurement environments. The welfare counterfactuals
(Section~\ref{sec:counterfactual}) are illustrative calculations that
combine the model's deployment logic with reduced-form price bounds.

\emph{Third}, we derive illustrative policy implications. The
minimum-bidder rule for convite procedures creates a corner solution
that forces cover-bidder deployment in 39.2\% of FL-present
tenders. Brazil's 2021 procurement reform (Lei~14.133) has since
abolished convite entirely, moving policy in a direction consistent
with this implication. We embed the detection ROC in a social
welfare function and show that the welfare-maximizing threshold
($1.6\times$IQR) coincides with the statistical optimum, though
the welfare calculations depend on the reduced-form price bounds.
The model also helps rationalize why participation-based screens
(AUC~$= 0.94$) outperform variance-based screens
(AUC~$= 0.79$) in this environment: under the coordinated regime,
bid-level statistics are uninformative.

\subsection{Related literature}

This paper connects to three strands. First, the structural auction
literature that estimates bid distributions to recover costs or detect
collusion \citep{guerre2000optimal, bajari2003deciding,
porter1993detection, pesendorfer2000study}. We contribute a mixture
model that separates cover bids from genuine bids using
participation-based labels rather than latent types, enabling direct
estimation of the cover-bid distribution without an EM algorithm.

Second, the proactive cartel screening literature
\citep{harrington2008detecting, imhof2019detecting,
huber2019machine}. Our structural model provides one mechanism-based
interpretation of the documented performance gap between
participation-based and variance-based screens in this setting.

Third, the procurement design literature
\citep{szymanski2006impact, decarolis2018buyers}. The model's corner solution
(Proposition~\ref{prop:optimal_m}) suggests that minimum-bidder
rules may create institutional demand for cover bidders in settings
where they bind.

The remainder of this paper is organized as follows.
Section~\ref{sec:structural} develops the structural model.
Section~\ref{sec:data} describes the data and FL classification.
Section~\ref{sec:empirical} details the estimation strategy.
Section~\ref{sec:results} presents the structural estimates,
specification tests, and detection performance comparison.
Section~\ref{sec:counterfactual} derives welfare counterfactuals.
Section~\ref{sec:robustness} reports robustness checks.
Section~\ref{sec:conclusion} concludes.
